\documentclass[11pt,letterpaper]{article}
\usepackage{ad42010-suite}
\begin{document}
\DocFront{01}{Architecture Description Fundamentals}{Conceptual model, terminology, and disciplined use}{A concept-first guide to describing the architecture of an entity of interest, preserving 2022 terminology and separating the standard's requirements from operational advice.}
\InterpretationKey
\section{Purpose and boundaries}
\begin{Requirement}
\FlagReq{} \textbf{Scope of the standard.} ISO/IEC/IEEE 42010:2022 addresses architecture \emph{descriptions}, not quality or correctness requirements for the architecture or the entity itself. Its conformance provisions apply separately to an architecture description (AD), architecture description framework (ADF), architecture description language (ADL), architecture viewpoint and model kind. \Source{Clauses 1 and 4}
\end{Requirement}
\begin{Guidance}Treat the AD as an independently versioned work product: it explains the architecture from selected viewpoints to interested stakeholders. A high-quality set of drawings alone is not automatically a conforming AD. Neither is a document guaranteed to describe a safe, performant or correct design just because it meets documentation requirements.\end{Guidance}
\section{Core vocabulary}
\small
\begin{longtable}{p{0.225\linewidth}p{0.51\linewidth}p{0.19\linewidth}}
\toprule\textbf{2022 term} & \textbf{Working interpretation} & \textbf{Reference}\\\midrule\endhead
Entity of interest (EoI)& Subject whose architecture is described; may be a solution, enterprise, service, process or other entity. & 3.12, 5.2.1\\
Architecture& Fundamental concepts or properties of the entity in its environment, including governing principles. & 3.2, 5.2.2\\
Architecture description (AD)& Tangible work product expressing an architecture; distinct from the architecture itself. & 3.3, 5.2.2\\
AD element& Identified or named part of the AD, including conceptual and view-related occurrences. & 3.4, 5.2.9\\
Stakeholder& Role, person, class or organization having an interest, right or claim in the EoI. & 3.17, 5.2.3\\
Concern& Matter relevant to a stakeholder and the AD's intended purpose. & 3.10, 5.2.3\\
Stakeholder perspective& Way of considering the EoI, particularly in relation to concerns. & 3.18, 5.2.4\\
Aspect& Characteristic dimension of the EoI, e.g., structural, functional, behavioural or informational. & 3.9, 5.2.5\\
Architecture viewpoint& Conventions that frame concerns and govern the construction and interpretation of views. & 3.8, 5.2.7\\
Architecture view& Portion of the AD providing information to address concerns under a viewpoint. & 3.7, 5.2.7\\
View component& Separable portion of one or more views, model-based or non-model-based. & 3.19, 5.2.8\\
Model kind& Model category identified by its characteristics and modelling conventions. & 3.15, 5.2.8\\
Legend& Documentation of conventions used to interpret a view component. & 3.19 note, 5.2.8\\
Correspondence& Identified relation connecting AD elements, potentially across ADs. & 3.11, 5.2.11\\
Correspondence method& Rule or procedure for describing or checking correspondences and consistency. & 5.2.11, 6.9.3\\
View method& Construction, interpretation or analysis approach associated with a viewpoint. & 5.2.10, 8.3\\
ADF& Domain/community conventions, principles and practices for describing architectures. & 3.5, 5.4\\
ADL& A set of representations plus syntax, semantics and rules for expressing architecture. & 3.6, 5.4.4\\
\bottomrule
\end{longtable}
\normalsize
\section{Relationships among the concepts}
\begin{center}
\begin{tikzpicture}[node distance=8mm and 12mm, every node/.style={font=\footnotesize}, b/.style={draw=Navy,rounded corners,align=center,fill=Pale,minimum width=27mm,minimum height=9mm}, >=Latex]
\node[b] (st) {Stakeholder};
\node[b,right=of st] (co) {Concern};
\node[b,right=of co] (vp) {Architecture\\viewpoint};
\node[b,right=of vp] (vw) {Architecture\\view};
\node[b,below=of co] (pr) {Stakeholder\\perspective};
\node[b,below=of vp] (asp) {Aspect};
\node[b,below=of vw] (com) {View component};
\draw[->] (st) -- node[above]{holds} (co);
\draw[->] (vp) -- node[above]{governs} (vw);
\draw[->] (vp) -- node[above]{frames} (co);
\draw[->] (vw) -- node[right]{contains} (com);
\draw[->] (st) -- (pr);
\draw[->] (pr) -- (co);
\draw[->] (asp) -- (co);
\draw[->] (asp) -- (com);
\end{tikzpicture}
\end{center}
\begin{Guidance}The illustration is an original simplified conceptual aid, \emph{not} a reproduction or exhaustive metamodel of the standard. Cardinalities must not be inferred from drawing layout. Keep registries and identifiers for many-to-many associations: stakeholder--concern, concern--viewpoint, viewpoint--view, and view--component.\end{Guidance}
\section{The 2022 changes that matter}
\begin{itemize}
\item \Key{EoI over system of interest:} architecture descriptions may concern enterprises, services, businesses and other entities, not only systems.
\item \Key{ADF over architecture framework:} distinguishes architecture-description conventions from other frameworks for architecting and evaluating architectures.
\item \Key{Stakeholder perspectives and aspects:} first-class elements separating stakeholders' ways of thinking from characteristics of the EoI.
\item \Key{View components over architecture models:} a view may contain both model-based components and narrative or documentary components.
\item \Key{One viewpoint to multiple views:} preserve the explicit governing viewpoint--view relation and the concerns addressed by each view.
\item \Key{Cross-description correspondences:} supported as part of tracking related architectures and alternatives.
\end{itemize}
\Source{Foreword; 5.2.3--5.2.11}
\section{Conformance is not architecture endorsement}
\begin{Requirement}\FlagReq{} \textbf{Five conformance objects:} test ADs against Clause 6; ADFs against 7.1; ADLs against 7.2; viewpoint specifications against 8.1; and model-kind specifications against 8.2. The standard does not permit tailoring the obligations when claiming conformance. \Source{Clause 4}\end{Requirement}
\begin{Guidance}An internal readiness review should mark every applicable provision as demonstrated, gap, or nonconformance with rationale. Do not label a system ``42010 certified'', treat a tool as automatically compliant, or assert that documentation conformance establishes performance, security, safety or regulatory assurance.\end{Guidance}
\section{Normative and non-normative reading discipline}
\begin{tabularx}{\linewidth}{p{0.16\linewidth}X X}\toprule
\textbf{Category}&\textbf{Meaning in source}&\textbf{This suite's representation}\\\midrule
\FlagReq&Mandatory requirement from Clauses 6--8 or Clause 4&Paraphrased obligation with reference and evidence question.\\
\FlagRec&Recommendation&Promoted as a good practice only, never a conformance prerequisite by itself.\\
\FlagOpt&Permitted approach&One allowed option; not the only valid implementation.\\
\FlagImpl&Locally recommended engineering process&Original, optional controls; never mislabelled as a 42010 obligation.\\\bottomrule
\end{tabularx}
\begin{Worked}The standard calls for identifying concerns and connecting them to stakeholders (Clause 6.4). A local engineering team might implement this using tickets, YAML records, or a spreadsheet. The relationship is normative; the storage medium, naming convention and ticketing platform are not.\end{Worked}
\section{Architecture description through the life cycle}
\begin{enumerate}
\item \textbf{Frame intent}: identify EoI boundary and environment, scope, AD purpose and decisions to support.
\item \textbf{Elicit considerations}: identify stakeholders, perspectives, concerns and aspects without collapsing these distinct concepts.
\item \textbf{Select viewpoints}: choose conventions sufficient to frame all identified concerns; record versions and sources.
\item \textbf{Construct views}: make model-based or narrative view components following governing conventions, with appropriate legends.
\item \textbf{Reconcile}: document correspondences, known inconsistencies and methods used to assess consistency.
\item \textbf{Decide and maintain}: retain essential decisions, rationale, review history and change-impact evidence.
\end{enumerate}
\begin{Guidance}This sequence is a pragmatic workflow, not a process mandated by ISO/IEC/IEEE 42010. That standard does not prescribe an architecting process, tool, notation, format or medium. Life-cycle integration can be aligned separately to relevant systems and software engineering process standards.\end{Guidance}
\section{Worked EoI: a payment processing service}
\begin{Worked}
\textbf{EoI} \texttt{PAY-SVC-01}: an API-driven payment processing service. \textbf{Environment}: clients, identity provider, payment provider, monitoring service, cloud platform and regulatory context. \textbf{AD purpose}: explain architectural choices required for a security and operational readiness review.\par
\textbf{Stakeholders}: product owner \texttt{SH-01}, security engineer \texttt{SH-02}, platform operator \texttt{SH-03}.\par
\textbf{Perspectives}: business \texttt{SP-01}, operational \texttt{SP-02}, assurance \texttt{SP-03}.\par
\textbf{Concerns}: \texttt{CN-01} ``How is payment information protected?''; \texttt{CN-02} ``How do failures affect transaction completion?''; \texttt{CN-03} ``Who owns recovery?''\par
\textbf{Aspects}: informational \texttt{AS-01}, behavioural \texttt{AS-02}, structural \texttt{AS-03}.\par
\textbf{Viewpoints}: information-security \texttt{VP-SEC}, runtime \texttt{VP-RUN}, operational \texttt{VP-OPS}.\par
\textbf{Views}: trust-boundary \texttt{VW-SEC-01}, failure-flow \texttt{VW-RUN-01}, recovery \texttt{VW-OPS-01}.\par
\textbf{Components}: data-flow diagram \texttt{VC-01}, sequence model \texttt{VC-02}, runbook narrative \texttt{VC-03}.\par
\textbf{Correspondence}: \texttt{CO-01} requires each critical data-flow boundary to have a responsible control owner.\par
\textbf{Decision}: \texttt{ADR-001} selects tokenization at an external provider; includes alternatives and trade-offs.
\end{Worked}
\section{Practical questions for architecture reviewers}
\begin{itemize}
\item Can readers tell the EoI boundary from its environment and identify the intended purpose of this AD?
\item Can every concern be followed to its stakeholder, framing viewpoint and addressing view?
\item Can each aspect and stakeholder perspective be traced to a view?
\item Can a reader identify the governing conventions for each view component, including narratives?
\item Are known inconsistencies, unresolved issues, method violations and essential decisions visible?
\item Can someone explain why the AD was changed and how the change affects other views?
\end{itemize}
\LegalFooter
\end{document}
